% Adenda v2 a las conclusiones: cierra OE6 y H6 en el mismo vocabulario de la
% Tabla de estado de hipótesis, sin tocar
% sections/source-snapshot/mainmatter/07-conclusions.tex.
\subsection{Cierre de OE6, H6 y RQ6}
\label{subsec:conclusions-megajuego}

\textbf{RQ6.} Sí se certifica por bloque, dentro de los dominios que cada
resultado declara y sin que la cadena completa herede una garantía conjunta.
La segunda mitad de la pregunta admite una respuesta más débil de lo que su
enunciado sugiere: la conclusión cualitativa de separación del mundo pareado
se reproduce en CoppeliaSim mediante reejecución cinemática, pero esa
comprobación no revalida en ese motor las hipótesis dinámicas, de contacto ni
de tracción rueda--suelo del certificado de SP2. Reproducir un resultado
geométrico no es revalidar el certificado que lo respalda.

OE6 se cumple de forma delimitada. El juego de integración conecta SP1--SP3 con
cuatro resultados demostrados algebraicamente y verificados numéricamente
(Sección~\ref{subsec:megajuego}, Tabla~\ref{tab:megajuego-e2e-summary}), y la
reproducción en CoppeliaSim del juego de acoplamiento de E4
(Sección~\ref{subsec:coppelia}, Tabla~\ref{tab:coppelia-compact-narrow})
conserva el resultado cualitativo en un
motor distinto del que lo generó: el método directo colisiona en el mundo
pareado de paso estrecho y el juego distribuido no. Lo que OE6 no entrega es un
certificado de optimalidad global del sistema híbrido bajo revisiones discretas
arbitrarias; la Proposición~\ref{prop:pre-cross-partition-limit} da la razón
formal de esa ausencia.

\begin{table}[H]
  \centering
  \caption{Estado final de H6, en continuación de la Tabla~\ref{tab:conclusion-hypotheses}.}
  \label{tab:conclusion-hypothesis-h6}
  \scriptsize
  \setlength{\tabcolsep}{3pt}
  \begin{tabularx}{\textwidth}{@{}p{0.8cm}p{1.75cm}p{5.1cm}X@{}}
    \toprule
    ID & Estado & Evidencia principal & Límite \\
    \midrule
    H6 & Apoyo delimitado & Cuatro resultados con evidencia de tipo distinto: potencial por factores y equivalencia variacional, por demostración algebraica; separación de Bézier, por demostración exacta más extragradiente sobre 30 instancias; \emph{caging}, por KKT racionales con residual $1{,}4\times10^{-11}$; presupuesto de ejecución, por desigualdad condicional. La reproducción en CoppeliaSim conserva el resultado cualitativo de separación (Tabla~\ref{tab:coppelia-compact-narrow}). & Dominios concretos por resultado; solo dos de los cuatro tienen verificación numérica. La reproducción es cinemática y de una corrida. \\
    \bottomrule
  \end{tabularx}
  \viuownsource
\end{table}

Los dominios son los que se demuestran: la separación Bézier, a dos coaliciones
con una familia geométrica fija de curvas; el \emph{caging} racional, a una
carga circular con cuatro contactos cardinales; la reproducción en CoppeliaSim,
a geometría y reejecución cinemática, sin contacto ni hardware. El trabajo
futuro más directo es la campaña congelada del
Anexo~\ref{app:megajuego-full-detail} (misma planta, mismo objetivo, control
central de referencia, variando agentes, contacto, batería, fricción, retardo y
presupuesto de cálculo) junto con extender CoppeliaSim a contacto y dinámica de
ruedas antes de un piloto de hardware.

Una política aprendida no sustituye a esta arquitectura porque entrega una
acción, no un predicado verificable con dominio declarado, y sin predicado no
hay columna que rellenar en la regla de composición de la
Tabla~\ref{tab:pre-cross-certificates}. Sí encaja como generador de
propuestas dentro del filtro: la política propone y el certificado de
\emph{wrench} admite o se abstiene, lo que conserva el contrato y es la vía de
extensión más inmediata para la familia de métodos que SP1 ya compara.

\subsection{Recomendaciones por dimensión}
\label{subsec:conclusions-recommendations-dimensions}

Las recomendaciones enunciadas hasta aquí son de implementación. Tres
dimensiones adicionales condicionan si ese trabajo merece hacerse, y ninguna se
resuelve con más simulación.

En lo teórico, la garantía extremo a extremo que hoy no existe requiere dos
piezas concretas que la Sección~\ref{subsec:cross-discussion} ya
identifica: una función de Lyapunov común, o un argumento de tiempo de
permanencia, para la conmutación entre modos de contacto, y una prueba de
compatibilidad entre los tres modelos de planta que hoy solo se compone de
forma funcional. Sin ellas, la composición de certificados seguirá siendo una
regla de interfaz y no un teorema.

En lo económico, la decisión relevante no es el coste de la flota frente al de
un manipulador único, sino el coste de cómputo por misión: el certificado de
\emph{wrench} resuelve un QP por admisión y la reserva espacio--temporal se
reevalúa en cada revisión, de modo que el criterio de viabilidad es si ese coste
cabe en el presupuesto de ejecución del
Teorema~\ref{thm:megajuego-execution-budget} con el \emph{hardware} de a
bordo previsto, no en una estación remota. Medir ese coste en la plataforma
objetivo debe preceder a cualquier compra.

En lo social, la abstención es la propiedad con consecuencia directa sobre el
operario. Un sistema que declara «no puedo certificar este transporte» y se
detiene traslada la decisión a una persona, y esa persona necesita saber qué
condición falló; recomendar el despliegue exige, por tanto, exponer la causa de
la abstención en la interfaz de planta, no solo el estado. La manipulación
pesada asistida es precisamente el caso donde un falso positivo de admisión
tiene coste físico sobre quien está cerca de la carga.

\subsection{Condiciones bajo las que esta arquitectura no debe usarse}
\label{subsec:conclusions-do-not-use}

Los límites anteriores se declaran por resultado. Reunidos, delimitan el uso del
conjunto. La arquitectura no debe aplicarse con congestión alta: en el piloto
Industrial~2, tres de los cuatro escenarios entregaron cero con los cuatro
métodos, incluidas las referencias, y ese régimen no está caracterizado. Tampoco
con cargas no planares, contactos móviles durante el transporte o más de una
carga simultánea sujeta a un oráculo espacio--temporal, porque ninguno de los
tres casos entra en el dominio de los certificados de SP2 y SP3. Fuera del
modelo de pérdidas Bernoulli --- particiones persistentes, retardos no
acotados --- la mensajería vecinal no tiene garantía ensayada. Y ningún
resultado de esta memoria sostiene despliegue físico: el verificador registró
dos violaciones de barrera bajo el modelo declarado, lo que basta para exigir
una cota continua, no muestreada, antes de cualquier prueba con hardware.
